% Generated draft: requires mathematical and source review.
\documentclass[11pt]{article}
\usepackage[margin=25mm]{geometry}
\usepackage{fontspec}
\setmainfont{Latin Modern Roman}
\usepackage{amsmath,amssymb,mathtools,unicode-math}
\setmathfont{Latin Modern Math}
\usepackage{enumitem,needspace,xurl,hyperref}
\setlength{\parindent}{0pt}
\setlength{\parskip}{5pt}
\setlength{\emergencystretch}{3em}
\newcommand{\N}{\mathbb{N}}
\newcommand{\Z}{\mathbb{Z}}
\newcommand{\Q}{\mathbb{Q}}
\newcommand{\R}{\mathbb{R}}
\newcommand{\C}{\mathbb{C}}
\newcommand{\F}{\mathbb{F}}
\newcommand{\A}{\mathbb{A}}
\newcommand{\PP}{\mathbb{P}}
\newcommand{\E}{\mathbb{E}}
\newcommand{\eps}{\varepsilon}
\newcommand{\dd}{\,\mathrm{d}}
\begin{document}
\section*{UnsolvedMath ID: 10400166}
\textbf{Source problem code:} catalogue entry
\subsection{Definitions and mathematical statement}
Let \(Z\) be a given three-dimensional TQFT in the source's conventions. A typical finite fusion-rule algebra has a distinguished finite basis \(\{x_i:i\in I\}\), a unit, duality data, and multiplication
\[
 x_i x_j=\sum_{k\in I}N_{ij}^{k}x_k,
 \qquad N_{ij}^{k}\in\Z_{\ge0},
\]
subject to the appropriate associativity and duality conditions. Associated \(6j\)-symbols encode changes of basis for iterated fusion and must satisfy the relevant coherence equations, including a pentagon identity. This is descriptive terminology: the exact coefficient system, normalization and extra symmetries in the source's class are not supplied.

Y. Kawahigashi's Problem 9.2 asks:
\begin{enumerate}
\item Can one determine whether the given \(Z\) arises from a fusion-rule algebra together with \(6j\)-symbols?
\item If it does, can one describe all fusion-rule algebras with \(6j\)-symbols that produce \(Z\)?
\end{enumerate}
A complete version must specify the construction \(D\mapsto Z_D\) from such data and the meaning of equivalence \(Z_D\cong Z\). The first task is to characterize whether \(Z\) belongs to that construction's image; the second is to describe its full collection of preimages, with whatever equivalence on fusion data the source intends. The question does not assert uniqueness. Nor does it provide a finite encoding that would turn ``determine'' into a formal algorithmic-decidability claim. The TQFT category, orientation, unitary or semisimple requirements, extension level, and permissible fusion data must be recovered before choosing a particular state-sum framework.
\subsection{Short English statement}
Determine whether a given three-dimensional TQFT comes from fusion rules and coherent \(6j\)-symbols, then classify all such input data producing it. The source does not define the construction or the equivalences under which theories and fusion data are compared.
\subsection{Sources}
\begin{itemize}
\item Local UnsolvedMath registry, permanent ID \(10400166\), exact source code \texttt{catalogue entry}, Problem 9.2, attributed to Y. Kawahigashi. The supplied title is truncated after ``(Y.''
\item Ohtsuki, \emph{Problems on invariants of knots and 3-manifolds} (2002), Problem 9.2, PDF page 138: \url{https://msp.org/gtm/2002/04/gtm-2002-04-024s.pdf}. Bold printed headings distinguish 231 numbered items from surrounding references.
\item Extraction pdf-font-structure, accessed 2026-07-29. Source evidence NEEDS\_REVIEW presents an open problem and requires a current primary-literature audit. Supplied research classification OPEN-TRIAGE identifies this fusion-rule-algebra question without a resolution or additional references.
\item Dataset attribution: Ulam AI / UnsolvedMath Contributors, CC BY 4.0. Source rights NEEDS\_REVIEW: publicly accessible author/publisher copy, redistribution terms requiring review.
\end{itemize}
\end{document}
